% ── chemistry-handout-preamble.tex ── 化学竞赛讲义 LaTeX 导言区（v4 ❄️）
% 2026-06-27 全面优化版

\documentclass[a4paper,11pt]{ctexart}

% ── 编码 ──
\usepackage[T1]{fontenc}

% ── 化学 ──
\usepackage[version=4]{mhchem}

% ── 数学 ──
\usepackage{amsmath,amssymb}
\usepackage{breqn}           % dmath 环境自动换行长公式
\usepackage{mathtools}       % 数学扩展

% ── 图片 ──
\usepackage{graphicx}

% ── 表格（跨页+自动列宽）──
\usepackage{booktabs}
\usepackage{xltabular}       % longtable + tabularx 合并
\usepackage{array}
\newcolumntype{L}{>{\raggedright\arraybackslash}X}
\newcolumntype{C}{>{\centering\arraybackslash}X}
\newcolumntype{R}{>{\raggedleft\arraybackslash}X}

% ── 浮动体（[H] 强制定位）──
\usepackage{float}

% ── 页面 ──
\usepackage[includeheadfoot,margin=2.2cm,headheight=14pt]{geometry}
\setlength{\textwidth}{16.6cm}

% ── 超链接(无红框) ──
\usepackage[colorlinks=true,linkcolor=black,citecolor=black,urlcolor=black]{hyperref}

% ── 颜色 ──
\usepackage[dvipsnames,svgnames]{xcolor}
\definecolor{chapterblue}{RGB}{23,50,77}
\definecolor{insightbg}{RGB}{235,240,250}
\definecolor{warningbg}{RGB}{255,245,235}
\definecolor{tipbg}{RGB}{235,250,235}
\definecolor{speechbg}{RGB}{245,235,250}

% ── 页眉页脚（精致化）──
\usepackage{fancyhdr}
\pagestyle{fancy}
\fancyhf{}
\fancyhead[L]{\small\color{chapterblue}\leftmark}
\fancyhead[R]{\small\color{chapterblue}\thepage}
\renewcommand{\headrule}{{\color{chapterblue}\hrule height 0.8pt}}

% ── 章节编号与样式 ──
\setcounter{secnumdepth}{1}
\ctexset{
  section = {
    format = \Large\bfseries\heiti\color{chapterblue},
    aftername = \hspace{0.5em},
    beforeskip = 2ex,
    afterskip = 1ex,
  },
  subsection = {
    format = \large\bfseries\heiti,
    beforeskip = 1.5ex,
    afterskip = 0.5ex,
  },
}

% ── 教学提示框（带浅底色，无 tcolorbox，完全兼容 dvipdfmx）──
\newenvironment{insightbox}{%
  \par\smallskip
  \noindent{\small\bfseries\color{RoyalBlue}🧠 认知冲突}\par
  \noindent{\small\color{black!80}}%
  \leavevmode\colorbox{insightbg}{\parbox{\dimexpr\textwidth-2\fboxsep\relax}{%
  \ignorespaces}}{%
  }\par\smallskip}

\newenvironment{warningbox}{%
  \par\smallskip
  \noindent{\small\bfseries\color{BrickRed}⚠️ 易错点}\par
  \noindent{\small\color{black!80}}%
  \leavevmode\colorbox{warningbg}{\parbox{\dimexpr\textwidth-2\fboxsep\relax}{%
  \ignorespaces}}{%
  }\par\smallskip}

\newenvironment{tipbox}{%
  \par\smallskip
  \noindent{\small\bfseries\color{ForestGreen}💡 理解提示}\par
  \noindent{\small\color{black!80}}%
  \leavevmode\colorbox{tipbg}{\parbox{\dimexpr\textwidth-2\fboxsep\relax}{%
  \ignorespaces}}{%
  }\par\smallskip}

\newenvironment{speechbox}{%
  \par\smallskip
  \noindent{\small\bfseries\color{Plum}🗣️ 教师原话}\par
  \noindent{\small\color{black!80}}%
  \leavevmode\colorbox{speechbg}{\parbox{\dimexpr\textwidth-2\fboxsep\relax}{%
  \ignorespaces}}{%
  }\par\smallskip}

% ── 列表间距 ──
\usepackage{enumitem}
\setlist{nosep}

% ── 防止内容溢出 ──
\usepackage{microtype}       % 字符伸缩，改善断行质量
\setlength{\emergencystretch}{2em}
\setlength{\hfuzz}{2pt}
\setlength{\overfullrule}{0pt}
\sloppy

% ── 代码高亮（用于解答/注释）──
\usepackage{listings}
\lstset{
  basicstyle=\small\ttfamily,
  breaklines=true,
  frame=none,
}

% ── pandoc 兼容 ──
\providecommand{\tightlist}{\setlength{\itemsep}{0pt}\setlength{\parskip}{0pt}}
\newenvironment{Shaded}{}{}
\newenvironment{Highlighting}{}{}
\newcommand{\NormalTok}[1]{#1}
\newcommand{\CommentTok}[1]{#1}
\newcommand{\KeywordTok}[1]{#1}
\newcommand{\StringTok}[1]{#1}
\newcommand{\FunctionTok}[1]{#1}
\newcommand{\DataTypeTok}[1]{#1}
\newcommand{\NumberTok}[1]{#1}
\newcommand{\ControlFlowTok}[1]{#1}
\newcommand{\OperatorTok}[1]{#1}
\newcommand{\BaseNTok}[1]{#1}
\newcounter{none}

% ── 图片尺寸限制 ──
\AtBeginDocument{\setkeys{Gin}{width=\textwidth,keepaspectratio}}

% ── 参考/列表前手动间距（用于答案区）──
\newcommand{\answeritem}{\vspace{0.5em}}

% ── 防止图片跨节 ──
\usepackage{placeins}

\graphicspath{{C:/Temp/chem_media/}}
\title{配位化合物基础}
\date{}
\setlength{\tabcolsep}{3.5pt}
\small
\begin{document}
\maketitle
\tableofcontents
\newpage

\section{学生讲义：配位化合物基础（超级充实版·自学完整）}\label{ux5b66ux751fux8bb2ux4e49ux914dux4f4dux5316ux5408ux7269ux57faux7840ux8d85ux7ea7ux5145ux5b9eux7248ux81eaux5b66ux5b8cux6574}

\subsection{学习目标}\label{ux5b66ux4e60ux76eeux6807}

\begin{itemize}
\tightlist
\item
  ✅ 能说出配合物的组成（中心离子/原子、配体、配位原子、配位数、内界/外界），并用实验现象区分内界和外界
\item
  ✅ 能按 IUPAC 规则命名常见配合物（含配正离子、配负离子、中性配合物），并完成''名称↔化学式''的双向翻译
\item
  ✅ 能用价键理论（杂化轨道）解释配合物的空间构型和磁性，区分内轨/外轨型
\item
  ✅ 能用 EAN 18 电子规则判断羰基配合物的稳定性，理解奇电子配合物的稳定方式
\item
  ✅ 能用晶体场理论解释 d 轨道分裂机制，计算八面体/四面体/正方形场的 CFSE
\item
  ✅ 能根据 Δ 与 P 的相对大小判断高/低自旋，从磁矩数据反推 d 电子排布
\item
  ✅ 能用光谱化学序列比较配体场强，理解 σ 给体/π 给体/π 受体的分类框架
\item
  ✅ 能区分结构异构、几何异构和光学异构，用固定-变换法枚举异构体并判断手性
\item
  ✅ 能综合运用以上知识解决竞赛级配合物命名、构型判断、磁性分析和异构体枚举题
\end{itemize}

\subsection{一、核心概念}\label{ux4e00ux6838ux5fc3ux6982ux5ff5}

\textbf{关联知识点}：配合物 · 配位键 · 配位数

\subsubsection{1.1 配位键不是''特殊键''------它是 Lewis 酸碱对的自然延伸}\label{ux914dux4f4dux952eux4e0dux662fux7279ux6b8aux952eux5b83ux662f-lewis-ux9178ux78b1ux5bf9ux7684ux81eaux7136ux5ef6ux4f38}

从高中化学我们知道，NH3 有孤对电子，Cu2+ 有空轨道------当 Cu2+ 遇到 NH3，孤对电子进入空轨道，形成 Cu←NH3 配位键。这个过程的本质是：\textbf{Lewis 碱（配体）提供电子对给 Lewis 酸（中心离子/原子）}。

换句话说，配位键不是一种''特殊的''化学键------它就是 Lewis 酸碱反应的产物，只不过成键的电子对\textbf{来自同一个原子}（配体），而不是两个原子各出一个。但一旦形成，配位键与普通共价键没有本质区别。

\begin{quote}
🧠 \textbf{认知冲突}：「配位键形成后，还''特殊''吗？」------配位键与普通共价键的差异只在形成阶段（电子对的来源不同），一旦形成，\textbf{两者完全相同}。例如 \([Co(NH_3)_6]^{3+}\) 中的所有 Co---N 键都是等价的------无法区分哪个是''配位键''哪个是''普通共价键''。
\end{quote}

这就是配合物的本质：\textbf{中心原子（Lewis 酸）与配体（Lewis 碱）通过配位键结合形成的复杂化合物}。

\subsubsection{1.2 配合物的组成}\label{ux914dux5408ux7269ux7684ux7ec4ux6210}

配合物的结构可以写作：

\begin{Shaded}
\begin{Highlighting}[]
\NormalTok{配合物 = [内界] 外界}
\NormalTok{          ↑       ↑}
\NormalTok{      配位单元   离子键结合}
\end{Highlighting}
\end{Shaded}

{\def\LTcaptype{none} % do not increment counter
\begin{longtable}[]{@{}
  >{\raggedright\arraybackslash}p{(\linewidth - 4\tabcolsep) * \real{0.3333}}
  >{\raggedright\arraybackslash}p{(\linewidth - 4\tabcolsep) * \real{0.3333}}
  >{\raggedright\arraybackslash}p{(\linewidth - 4\tabcolsep) * \real{0.3333}}@{}}
\toprule\noalign{}
\begin{minipage}[b]{\linewidth}\raggedright
组成部分
\end{minipage} & \begin{minipage}[b]{\linewidth}\raggedright
定义
\end{minipage} & \begin{minipage}[b]{\linewidth}\raggedright
示例（{[}Cu(NH3)4{]}SO4）
\end{minipage} \\
\midrule\noalign{}
\endhead
\bottomrule\noalign{}
\endlastfoot
\textbf{中心离子/原子} & 位于配离子中心，接受孤对电子的原子/离子 & Cu2+ \\
\textbf{配体 (Ligand)} & 提供孤对电子的离子或分子 & NH3（4 个） \\
\textbf{配位原子} & 配体中直接提供孤对电子的原子 & N（NH3 中的 N） \\
\textbf{配位数} & 与中心离子直接键合的配位原子总数 & 4 \\
\textbf{内界} & 配位单元（配离子或中性分子），书写在方括号内 & {[}Cu(NH3)4{]}2+ \\
\textbf{外界} & 内界以外的部分，与内界以离子键结合 & SO42- \\
\end{longtable}
}

\begin{quote}
🧠 \textbf{认知冲突}：「配合物一定有外界吗？」------不一定！中性配合物如 \([Ni(CO)_4]\)、\([Fe(CO)_5]\) \textbf{只有内界，没有外界}。它们整体呈电中性，不需要外界离子来平衡电荷。
⚠️ \textbf{易错点}：配位数 ≠ 配体数！单齿配体时二者相等；多齿配体时\textbf{配位数 = 各配体齿数之和}。例如 \([Pt(en)_2]^{2+}\) 配体数为 2，但每个 en 是双齿配体，因此配位数 = 2 × 2 = 4。
\end{quote}

\textbf{判断内界/外界的关键实验}：

\textbf{例}：\([Cu(NH_3)_4]SO_4\) 溶液中：
- 加 NaOH → \textbf{无} Cu(OH)2 沉淀（几乎没有游离 Cu2+------Cu2+ 在内界被 NH3 牢固配位）
- 加 BaCl2 → \textbf{有} BaSO4 白色沉淀（SO42- 在外界，自由存在）
- 说明：溶液中主要存在 \([Cu(NH_3)_4]^{2+}\) 和 SO42-

\subsubsection{1.3 常见中心离子与配体}\label{ux5e38ux89c1ux4e2dux5fc3ux79bbux5b50ux4e0eux914dux4f53}

\paragraph{中心离子}\label{ux4e2dux5fc3ux79bbux5b50}

\begin{itemize}
\tightlist
\item
  \textbf{过渡金属离子}：Fe2+/Fe3+、Co2+/Co3+、Ni2+、Cu2+/Cu+、Zn2+、Cr3+、Mn2+ 等（最常见，有 d 轨道）
\item
  \textbf{高氧化态主族金属/非金属}：Al3+（\([AlF_6]^{3-}\)）、Si(IV)（\([SiF_6]^{2-}\)）
\item
  \textbf{中性原子}：Ni(0) 在 \([Ni(CO)_4]\) 中
\end{itemize}

\begin{quote}
🗣️ \textbf{课堂原话}：``过渡金属能做中心离子，是因为它们有\textbf{空的价层轨道}接受孤对电子------不全是 d 轨道，4s 和 4p（对第一过渡系）也参与。''
\end{quote}

\paragraph{常见配体}\label{ux5e38ux89c1ux914dux4f53}

{\def\LTcaptype{none} % do not increment counter
\begin{longtable}[]{@{}
  >{\raggedright\arraybackslash}p{(\linewidth - 4\tabcolsep) * \real{0.3077}}
  >{\raggedright\arraybackslash}p{(\linewidth - 4\tabcolsep) * \real{0.3077}}
  >{\centering\arraybackslash}p{(\linewidth - 4\tabcolsep) * \real{0.3846}}@{}}
\toprule\noalign{}
\begin{minipage}[b]{\linewidth}\raggedright
配体类型
\end{minipage} & \begin{minipage}[b]{\linewidth}\raggedright
实例
\end{minipage} & \begin{minipage}[b]{\linewidth}\centering
配位原子
\end{minipage} \\
\midrule\noalign{}
\endhead
\bottomrule\noalign{}
\endlastfoot
单齿（1 个配位原子） & NH3、H2O、CO、CN-、Cl-、F-、OH-、NO2-、SCN- & N、O、C、卤素、S \\
双齿（2 个配位原子） & en（乙二胺 H2N--CH2--CH2--NH2）、ox（草酸根 --OOC--COO-）、gly-（氨基乙酸 H2N--CH2--COO-） & N、O \\
多齿（≥3 个配位原子） & EDTA（6齿）、二乙烯三胺（3齿）、冠醚、穴醚 & N、O \\
\end{longtable}
}

\begin{quote}
⚠️ \textbf{易错点}：两可配体（ambidentate ligand）的配位原子需要根据情况判断。NO2- 可通过 N 配位（硝基）也可通过 O 配位（亚硝酸根）；SCN- 可通过 S 配位（硫氰酸根）也可通过 N 配位（异硫氰酸根）。这会导致\textbf{键合异构}。
\end{quote}

\subsubsection{1.4 配位数}\label{ux914dux4f4dux6570}

常见配位数：\textbf{2、4、6}（少数 3、5、8）

\textbf{影响配位数的因素}：

{\def\LTcaptype{none} % do not increment counter
\begin{longtable}[]{@{}
  >{\raggedright\arraybackslash}p{(\linewidth - 4\tabcolsep) * \real{0.3333}}
  >{\raggedright\arraybackslash}p{(\linewidth - 4\tabcolsep) * \real{0.3333}}
  >{\raggedright\arraybackslash}p{(\linewidth - 4\tabcolsep) * \real{0.3333}}@{}}
\toprule\noalign{}
\begin{minipage}[b]{\linewidth}\raggedright
因素
\end{minipage} & \begin{minipage}[b]{\linewidth}\raggedright
趋势
\end{minipage} & \begin{minipage}[b]{\linewidth}\raggedright
解释
\end{minipage} \\
\midrule\noalign{}
\endhead
\bottomrule\noalign{}
\endlastfoot
中心离子氧化数 & 氧化数越高 → 配位数越大 & Co3+ 配位数通常为 6，Co2+ 可为 4 或 6 \\
中心离子半径 & 半径越大 → 配位数一般越大 & 但过大时配体间排斥增大→配位数反而下降（Hg2+ 配位数 4） \\
配体体积 & 配体越大 → 配位数越小 & Al3+ 与 F-（小）形成 \([AlF_6]^{3-}\)（配位数 6），与 Cl-（大）形成 \([AlCl_4]^-\)（配位数 4） \\
配体电荷 & 配体负电荷越高 → 配位数越小 & 配体间静电排斥增大 \\
\end{longtable}
}

\begin{quote}
📝 \textbf{练一练}：判断 \([Pt(en)_2]^{2+}\) 中 Pt2+ 的配位数（en 为乙二胺，双齿配体）。
\end{quote}

\textbf{答案}：Pt2+ 配位数 = 2 个 en × 每个 en 提供 2 个配位原子 = \textbf{4}

\subsubsection{1.5 配合物的命名------``名称↔化学式''双向翻译}\label{ux914dux5408ux7269ux7684ux547dux540dux540dux79f0ux5316ux5b66ux5f0fux53ccux5411ux7ffbux8bd1}

\paragraph{命名顺序}\label{ux547dux540dux987aux5e8f}

\begin{verbatim}
配位数(汉字) → 配体名称 → "合" → 中心离子名称 → (氧化数罗马数字)
\end{verbatim}

\paragraph{配体顺序规则}\label{ux914dux4f53ux987aux5e8fux89c4ux5219}

\begin{enumerate}
\def\labelenumi{\arabic{enumi}.}
\tightlist
\item
  \textbf{先负离子，后中性分子}：\([Co(NO_2)_4(NH_3)_2]^-\) → 四硝基·二氨合钴(Ⅲ)离子
\item
  \textbf{先无机，后有机}：\([CoCl(SCN)(en)_2]^+\) → 氯·硫氰酸根·二(乙二胺)合钴(Ⅲ)离子
\item
  \textbf{同类配体按配位原子元素符号字母顺序排列}
\item
  \textbf{配位原子相同，含原子数少的配体排前}
\item
  多种配体之间用 \textbf{·}（中圆点）隔开
\end{enumerate}

\begin{quote}
🗣️ \textbf{课堂原话}：``命名不是死记硬背------它是『名称↔结构』的双向翻译能力。看到 K3{[}Fe(CN)6{]} 能写出六氰合铁(III)酸钾，反过来也要能写。''
\end{quote}

\paragraph{配合物命名速查表}\label{ux914dux5408ux7269ux547dux540dux901fux67e5ux8868}

{\def\LTcaptype{none} % do not increment counter
\begin{longtable}[]{@{}
  >{\raggedright\arraybackslash}p{(\linewidth - 4\tabcolsep) * \real{0.3333}}
  >{\raggedright\arraybackslash}p{(\linewidth - 4\tabcolsep) * \real{0.3333}}
  >{\raggedright\arraybackslash}p{(\linewidth - 4\tabcolsep) * \real{0.3333}}@{}}
\toprule\noalign{}
\begin{minipage}[b]{\linewidth}\raggedright
类型
\end{minipage} & \begin{minipage}[b]{\linewidth}\raggedright
命名规则
\end{minipage} & \begin{minipage}[b]{\linewidth}\raggedright
示例
\end{minipage} \\
\midrule\noalign{}
\endhead
\bottomrule\noalign{}
\endlastfoot
\textbf{含配正离子}（外界为简单阴离子） & \textbf{某化某} & \([Co(NH_3)_6]Cl_3\) 三氯化六氨合钴(Ⅲ) \\
\textbf{含配正离子}（外界为复杂阴离子） & \textbf{某酸某} & \([Cu(NH_3)_4]SO_4\) 硫酸四氨合铜(Ⅱ) \\
\textbf{含配负离子}（外界为 H+） & \textbf{某酸} & \(H_2[SiF_6]\) 六氟合硅(Ⅳ)酸 \\
\textbf{含配负离子}（外界为其他阳离子） & \textbf{某酸某} & \(K_4[Fe(CN)_6]\) 六氰合铁(Ⅱ)酸钾 \\
\textbf{无外界（中性配合物）} & 配体名称 + 合 + 中心原子 & \([Fe(CO)_5]\) 五羰基合铁(0) \\
\textbf{配阴离子为内界} & 词尾用''酸'' & \(K_3[Fe(CN)_6]\) → 六氰合铁(Ⅲ)\textbf{酸钾} \\
\textbf{配阳离子为内界} & 词尾用''化'' & \([Co(NH_3)_6]Cl_3\) → 三氯化六氨合钴(Ⅲ) \\
\end{longtable}
}

\begin{quote}
📝 \textbf{练一练}：命名下列配合物，并指出中心离子氧化数和配位数：
(1) \([Co(NH_3)_5Cl]Cl_2\)
(2) \(K_2[Zn(OH)_4]\)
(3) \([CrCl_2(H_2O)_4]Cl\)
\end{quote}

\textbf{参考答案}：
(1) 二氯化一氯·五氨合钴(Ⅲ)------Co3+，配位数 6
(2) 四羟基合锌(Ⅱ)酸钾------Zn2+，配位数 4
(3) 氯化二氯·四水合铬(Ⅲ)------Cr3+，配位数 6

相关 KP：配合物 · 配位键 · 配位数

\subsection{二、规律与方法}\label{ux4e8cux89c4ux5f8bux4e0eux65b9ux6cd5}

\subsubsection{2.1 价键理论（VB Theory）------从杂化轨道看配合物构型}\label{ux4ef7ux952eux7406ux8bbavb-theoryux4eceux6742ux5316ux8f68ux9053ux770bux914dux5408ux7269ux6784ux578b}

价键理论是最早成功解释配合物空间构型和磁性的理论。它的核心逻辑链非常简洁：\textbf{中心离子提供空轨道 → 空轨道杂化 → 杂化轨道接受配体孤对电子 → 杂化方式决定空间构型。}

\paragraph{要点}\label{ux8981ux70b9}

\begin{enumerate}
\def\labelenumi{\arabic{enumi}.}
\tightlist
\item
  中心离子有空的价层轨道，配体提供孤对电子
\item
  中心离子的空轨道先\textbf{杂化}，再接受配体孤对电子形成配位键
\item
  杂化类型决定配合物的\textbf{空间构型}
\end{enumerate}

\paragraph{杂化类型与空间构型}\label{ux6742ux5316ux7c7bux578bux4e0eux7a7aux95f4ux6784ux578b}

{\def\LTcaptype{none} % do not increment counter
\begin{longtable}[]{@{}
  >{\centering\arraybackslash}p{(\linewidth - 6\tabcolsep) * \real{0.2778}}
  >{\centering\arraybackslash}p{(\linewidth - 6\tabcolsep) * \real{0.2778}}
  >{\raggedright\arraybackslash}p{(\linewidth - 6\tabcolsep) * \real{0.2222}}
  >{\raggedright\arraybackslash}p{(\linewidth - 6\tabcolsep) * \real{0.2222}}@{}}
\toprule\noalign{}
\begin{minipage}[b]{\linewidth}\centering
配位数
\end{minipage} & \begin{minipage}[b]{\linewidth}\centering
杂化类型
\end{minipage} & \begin{minipage}[b]{\linewidth}\raggedright
空间构型
\end{minipage} & \begin{minipage}[b]{\linewidth}\raggedright
典型示例
\end{minipage} \\
\midrule\noalign{}
\endhead
\bottomrule\noalign{}
\endlastfoot
2 & sp & 直线形 & \([Ag(NH_3)_2]^+\)、\([Cu(NH_3)_2]^+\) \\
3 & sp2 & 平面三角形 & \([CuCl_3]^{2-}\)、\([HgI_3]^-\) \\
4 & sp3 & 正四面体 & \([ZnCl_4]^{2-}\)、\([Ni(NH_3)_4]^{2+}\) \\
4 & dsp2 & 平面正方形 & \([Ni(CN)_4]^{2-}\)、\([PtCl_4]^{2-}\)、\([Cu(NH_3)_4]^{2+}\) \\
5 & dsp3 & 三角双锥 & \([Fe(CO)_5]\)、\([Ni(CN)_5]^{3-}\) \\
6 & sp3d2 & 正八面体（外轨） & \([FeF_6]^{3-}\)、\([Ti(H_2O)_6]^{3+}\) \\
6 & d2sp3 & 正八面体（内轨） & \([Fe(CN)_6]^{3-}\)、\([Co(NH_3)_6]^{3+}\) \\
\end{longtable}
}

\begin{quote}
⚠️ \textbf{关键对比}：配位数相同时，杂化类型可以不同 → 空间构型不同。这是 VBT 的核心魅力------同样是 4 配位，Ni2+ 与 NH3 形成四面体（sp3），与 CN- 形成平面正方形（dsp2）。
\end{quote}

\paragraph{内轨型 vs 外轨型配合物}\label{ux5185ux8f68ux578b-vs-ux5916ux8f68ux578bux914dux5408ux7269}

{\def\LTcaptype{none} % do not increment counter
\begin{longtable}[]{@{}
  >{\raggedright\arraybackslash}p{(\linewidth - 4\tabcolsep) * \real{0.3333}}
  >{\raggedright\arraybackslash}p{(\linewidth - 4\tabcolsep) * \real{0.3333}}
  >{\raggedright\arraybackslash}p{(\linewidth - 4\tabcolsep) * \real{0.3333}}@{}}
\toprule\noalign{}
\begin{minipage}[b]{\linewidth}\raggedright
特征
\end{minipage} & \begin{minipage}[b]{\linewidth}\raggedright
外轨型
\end{minipage} & \begin{minipage}[b]{\linewidth}\raggedright
内轨型
\end{minipage} \\
\midrule\noalign{}
\endhead
\bottomrule\noalign{}
\endlastfoot
\textbf{杂化方式} & 外层轨道（ns, np, nd）如 sp3d2 & 内层 (n-1)d + ns + np，如 d2sp3 \\
\textbf{d 电子排布} & 不变，单电子数不变 & \textbf{重排}（电子归并），单电子数减少 \\
\textbf{磁性} & 不变（顺磁/高自旋） & 可能变为反磁/低自旋 \\
\textbf{键的性质} & 离子性较强，共价性较弱 & 共价性较强 \\
\textbf{稳定性} & 稳定性较差 & 稳定性较好 \\
\textbf{形成条件} & 弱场配体（F-、H2O） & 强场配体（CN-、CO） \\
\end{longtable}
}

\begin{quote}
🧠 \textbf{认知冲突}：「内轨型 = 低自旋，外轨型 = 高自旋------两个理论说的是同一回事吗？」------方向对但不完全等价。内/外轨型来自\textbf{价键理论}（关注哪些轨道参与杂化），高/低自旋来自\textbf{晶体场理论}（关注 Δ 与 P 的竞争）。在大多数情况下它们对应，但严格来说这是两个不同理论对同一现象的不同表述。
\end{quote}

\textbf{判断方法}：通过磁矩测定 → 计算单电子数 n → 判断是否发生 d 电子重排。

\textbf{磁矩公式}：

\[
\mu = \sqrt{n(n+2)}\ (\text{B.M.})
\]

其中 n 为未成对电子数。\(\mu = 0\) → 反磁性；\(\mu > 0\) → 顺磁性。

\textbf{例}：\([CoF_6]^{3-}\) 磁矩 5.26 B.M. → n = 4 → d 电子未重排 → sp3d2 外轨型。
\([Co(CN)_6]^{3-}\) 磁矩 0 → n = 0 → d 电子归并 → d2sp3 内轨型。

\begin{quote}
🧠 \textbf{认知冲突}：「为什么 {[}Ni(CN)4{]}2- 是平面四边形，而 {[}NiCl4{]}2- 是四面体？」------Ni2+ 是 d⁸ 构型，VSEPR 预测四面体（4 对电子）。但 CN- 是强场配体→迫使 Ni2+ 的 d 电子重排归并→空出一个 3d 轨道→dsp2 杂化→平面四边形。Cl- 是弱场配体→不能迫使电子重排→sp3 杂化→四面体。\textbf{同一个金属离子，配体不同，不仅颜色和磁性不同------连形状都不同。}
\end{quote}

\paragraph{价键理论的局限性}\label{ux4ef7ux952eux7406ux8bbaux7684ux5c40ux9650ux6027}

VBT 是一个\textbf{解释工具}，而非\textbf{预测工具}。它无法定量预测：
- 配合物的颜色（d-d 跃迁能量）
- 光谱化学序列的来源（为什么 CO 是强场？）
- 分裂能 Δ 的具体数值

这些问题需要\textbf{晶体场理论（CFT）}来解决。

相关 KP：配合物 · 配位键 · 杂化轨道理论

\subsubsection{2.2 EAN 18 电子规则------金属有机配合物的''八隅体''}\label{ean-18-ux7535ux5b50ux89c4ux5219ux91d1ux5c5eux6709ux673aux914dux5408ux7269ux7684ux516bux9685ux4f53}

\textbf{有效原子序数（EAN）规则}：中心金属的价电子数 + 配体提供的 σ 电子数之和 = 18（即同周期稀有气体电子数，对应 ns2np⁶nd¹⁰ 满壳层）。

\begin{quote}
🗣️ \textbf{课堂原话}：``18 电子规则就像配合物世界的『八隅体规则』------金属原子利用全部 9 条价轨道（5 d + 1 s + 3 p）达到 18 电子稳定结构。''
\end{quote}

\paragraph{电子计数方法}\label{ux7535ux5b50ux8ba1ux6570ux65b9ux6cd5}

{\def\LTcaptype{none} % do not increment counter
\begin{longtable}[]{@{}
  >{\raggedright\arraybackslash}p{(\linewidth - 4\tabcolsep) * \real{0.3077}}
  >{\centering\arraybackslash}p{(\linewidth - 4\tabcolsep) * \real{0.3846}}
  >{\raggedright\arraybackslash}p{(\linewidth - 4\tabcolsep) * \real{0.3077}}@{}}
\toprule\noalign{}
\begin{minipage}[b]{\linewidth}\raggedright
配体类型
\end{minipage} & \begin{minipage}[b]{\linewidth}\centering
提供的电子数
\end{minipage} & \begin{minipage}[b]{\linewidth}\raggedright
示例
\end{minipage} \\
\midrule\noalign{}
\endhead
\bottomrule\noalign{}
\endlastfoot
经典单齿配体（NH3、CO、Cl-、H-、R-） & 2 & CO 提供 2e- \\
NO（一氧化氮） & 3 & NO = NO+ + e-，NO+ 提供 2e-（直线形配位） \\
桥联 CO（M---CO---M） & 每条桥键为 1 个金属贡献 1e- & \([Fe_2(CO)_9]\) 中每个 Fe 从桥联 CO 获 3e- \\
M---M 金属键 & 1 & 每条金属键为 1 个金属贡献 1e- \\
ηⁿ 配体（如 η⁵-C5H5-） & n（奇数时从金属取 1 凑偶数） & η⁵-C5H5 提供 6e- \\
配合物所带电荷 & 负电荷加电子，正电荷减电子 & \([Co(CO)_4]^-\) 中 Co- 相当于 Co 多 1e- \\
\end{longtable}
}

\paragraph{应用}\label{ux5e94ux7528}

\textbf{① 判断稳定性}：18 电子结构稳定，奇电子（17 电子）不稳定。

\textbf{奇电子配合物的稳定方式}：
1. 从还原剂夺得 1e- 成为负离子：\([M(CO)_n]^-\)
2. 与 H+ 或 X- 结合：\(HM(CO)_n\)、\(M(CO)_nX\)
3. 二聚化：\(Mn_2(CO)_{10}\)------Mn(0) 7e-，每个 Mn: \(5 \times CO + 1 \times Mn—Mn = 10 + 1 = 11e-\) 来自配体，加上 7 个价电子 → 18 ✅

\begin{quote}
⚠️ \textbf{易错点}：18 电子规则是\textbf{必要不充分条件}------满足 18 e- 不意味着配合物一定稳定（还要看位阻、配体电子性质）。\([V(CO)_6]\)（17e-）不形成二聚体，因为配位数变为 7 时空阻太大。
\end{quote}

\textbf{② 预测反应方向}：
\([Cr(CO)_6] + C_6H_6 \rightarrow [Cr(C_6H_6)(CO)_3] + 3CO\)
苯为 6e- 给予体，取代 3 个 CO（每个 2e-），电子总数不变，反应可进行。

\textbf{③ 估算 M---M 键数（簇合物分析）}：
如 \(Ir_4(CO)_{12}\) → 总电子数 \(4 \times 9 + 12 \times 2 = 60\)，平均每个 Ir 15e-，需 3 条 Ir---Ir 键 → 正四面体簇

\begin{quote}
📝 \textbf{练一练}：验证下列配合物是否符合 18 电子规则：
(1) \([Cr(CO)_6]\)
(2) \([Mn_2(CO)_{10]}\)（Mn 为 7 族）
(3) \([Fe(CO)_4H_2]\)（H- 为 2e- 给体）
\end{quote}

\textbf{参考答案}：
(1) Cr(0): 6e- + 6×CO=12 → \textbf{18 ✅}
(2) Mn(0): 7e-，每个 Mn: 5×CO端基=10 + 1×Mn---Mn键=1 → \textbf{18 ✅}
(3) Fe(II): 6e- + 4×CO=8 + 2×H-=4 → \textbf{18 ✅}

相关 KP：18电子规则 · 反馈π键

\subsubsection{2.3 晶体场理论（CFT）------为什么配合物有颜色和磁性？}\label{ux6676ux4f53ux573aux7406ux8bbacftux4e3aux4ec0ux4e48ux914dux5408ux7269ux6709ux989cux8272ux548cux78c1ux6027}

如果说价键理论回答了''配合物是什么形状''，那么晶体场理论回答了更深层的问题：\textbf{为什么配合物有颜色？为什么有的配合物顺磁、有的反磁？为什么配体不同会导致磁性反转？}

\begin{quote}
🧠 \textbf{认知冲突}：「{[}Fe(CN)6{]}⁴- 反磁性，{[}FeF6{]}⁴- 顺磁性------同一个 Fe2+！」------Fe2+ 是 d⁶，按洪特规则应有 4 个未成对电子。但 {[}Fe(CN)6{]}⁴- 是反磁性的（无未成对电子）！这正是晶体场理论要解释的核心现象。
\end{quote}

\paragraph{基本要点}\label{ux57faux672cux8981ux70b9}

\begin{enumerate}
\def\labelenumi{\arabic{enumi}.}
\tightlist
\item
  配体视为\textbf{点电荷}（或点偶极），与中心离子间为静电作用
\item
  中心离子的 d 轨道受配体负电场排斥作用 → \textbf{能级分裂}
\item
  d 电子重排 → 体系能量变化 → \textbf{晶体场稳定化能（CFSE）}
\end{enumerate}

\paragraph{八面体场中的 d 轨道分裂}\label{ux516bux9762ux4f53ux573aux4e2dux7684-d-ux8f68ux9053ux5206ux88c2}

\begin{figure}[H]\centering
\includegraphics[width=0.9\textwidth,keepaspectratio]{octahedral_splitting.png}
\caption{图1 五个 d 轨道的空间取向 — 与配体相对位置决定能级分裂方式}
\end{figure}

在正八面体场中，6 个配体沿 ±x、±y、±z 方向接近中心离子。五个简并的 d 轨道分裂为两组：

{\def\LTcaptype{none} % do not increment counter
\begin{longtable}[]{@{}
  >{\centering\arraybackslash}p{(\linewidth - 6\tabcolsep) * \real{0.2632}}
  >{\raggedright\arraybackslash}p{(\linewidth - 6\tabcolsep) * \real{0.2105}}
  >{\centering\arraybackslash}p{(\linewidth - 6\tabcolsep) * \real{0.2632}}
  >{\centering\arraybackslash}p{(\linewidth - 6\tabcolsep) * \real{0.2632}}@{}}
\toprule\noalign{}
\begin{minipage}[b]{\linewidth}\centering
轨道组
\end{minipage} & \begin{minipage}[b]{\linewidth}\raggedright
轨道
\end{minipage} & \begin{minipage}[b]{\linewidth}\centering
与配体方向
\end{minipage} & \begin{minipage}[b]{\linewidth}\centering
能量变化
\end{minipage} \\
\midrule\noalign{}
\endhead
\bottomrule\noalign{}
\endlastfoot
\(e_g\)（dγ） & \(d_{x^2-y^2, d_{z^2}}\) & \textbf{指向}配体 → 排斥大 & \(+6\ Dq\)（升高） \\
\(t_{2g}\)（dε） & \(d_{xy, d_{yz}, d_{xz}}\) & \textbf{指向}配体之间 → 排斥小 & \(-4\ Dq\)（降低） \\
\end{longtable}
}

分裂能 \(\Delta_o = E(e_g) - E(t_{2g}) = 10\ Dq\)

\textbf{能量守恒验证}：\(2 \times (+6) + 3 \times (-4) = 0\) ✓（重心不变原理）

\begin{quote}
🗣️ \textbf{课堂原话}：``想象 5 个 d 轨道是 5 个人站在操场上------配体像 6 个拿着水枪的同学从 ±x、±y、±z 方向喷水。站得正对着水枪方向的（e\_g）被喷得最惨，站得斜的（t2g）躲开了。所以 e\_g 能量升高，t2g 能量降低。''
\end{quote}

\begin{figure}[H]\centering
\includegraphics[width=0.9\textwidth,keepaspectratio]{octahedral_splitting.png}
\caption{图2 八面体场中 d 轨道分裂 — t2g 降低 4 Dq，e_g 升高 6 Dq，Δₒ = 10 Dq}
\end{figure}

\begin{figure}[H]\centering
\includegraphics[width=0.9\textwidth,keepaspectratio]{octahedral_splitting.png}
\caption{图2b 八面体场 d 轨道分裂示意图 — 左：六配位八面体场中六个配体沿 ±x、±y、±z 方向接近；右：d 轨道分裂为 t2g（低能）和 eg（高能），Δₒ = 10 Dq，重心守恒}
\end{figure}

\paragraph{其他配位场中的分裂}\label{ux5176ux4ed6ux914dux4f4dux573aux4e2dux7684ux5206ux88c2}

{\def\LTcaptype{none} % do not increment counter
\begin{longtable}[]{@{}
  >{\centering\arraybackslash}p{(\linewidth - 6\tabcolsep) * \real{0.2632}}
  >{\centering\arraybackslash}p{(\linewidth - 6\tabcolsep) * \real{0.2632}}
  >{\raggedright\arraybackslash}p{(\linewidth - 6\tabcolsep) * \real{0.2105}}
  >{\centering\arraybackslash}p{(\linewidth - 6\tabcolsep) * \real{0.2632}}@{}}
\toprule\noalign{}
\begin{minipage}[b]{\linewidth}\centering
配位场
\end{minipage} & \begin{minipage}[b]{\linewidth}\centering
分裂能
\end{minipage} & \begin{minipage}[b]{\linewidth}\raggedright
轨道能级顺序（从高到低）
\end{minipage} & \begin{minipage}[b]{\linewidth}\centering
图示
\end{minipage} \\
\midrule\noalign{}
\endhead
\bottomrule\noalign{}
\endlastfoot
八面体（Oh） & \(\Delta_o = 10\ Dq\) & \(d_{x^2-y^2 = d_{z^2} > d_{xy} = d_{xz} = d_{yz}}\) & \(e_g \uparrow 6\ Dq\)，\(t_{2g} \downarrow 4\ Dq\) \\
四面体（Td） & \(\Delta_t = \frac{4}{9}\Delta_o \approx 4.45\ Dq\) & \(d_{xy = d_{xz} = d_{yz} > d_{x^2-y^2} = d_{z^2}}\) & \(t_2 \uparrow\)，\(e \downarrow\)（与 Oh \textbf{相反}） \\
平面正方形（D4h） & \(\Delta_s \approx 17.42\ Dq\) & \(d_{x^2-y^2 > d_{xy} > d_{z^2} > d_{xz} = d_{yz}}\) & dₓ2₋ᵧ2 最高，dₓ2=dᵧ2 最低 \\
\end{longtable}
}

\begin{quote}
⚠️ \textbf{易错点}：四面体场的轨道符号是 e 和 t2（\textbf{不是} e\_g 和 t2g）------因为四面体场没有对称中心（g 表示中心对称）。四面体场的分裂能约为八面体场的 4/9，这是因为：① 配位数少（4 vs 6）；② 配体与 d 轨道的重叠程度小。
\end{quote}

\begin{figure}[H]\centering
\includegraphics[width=0.9\textwidth,keepaspectratio]{tetrahedral_vs_octahedral.png}
\caption{图3 四面体场中 d 轨道分裂 — 与八面体场能级顺序相反（Δₜ 约为 Δₒ 的 4/9）}
\end{figure}

\begin{figure}[H]\centering
\includegraphics[width=0.9\textwidth,keepaspectratio]{tetrahedral_vs_octahedral.png}
\caption{图3b 不同几何场中 d 轨道分裂对比 — 左：Oh 场 t2g↓/eg↑（Δₒ）；中：Td 场 e↓/t2↑（Δₜ≈4/9Δₒ）；右：D4h 场四层分裂（Δₛₚ > Δₒ）。注意 Td 场的分裂顺序与 Oh 相反}
\end{figure}

\paragraph{空白互动表：Oh/Td/D4h 分裂参数对比}\label{ux7a7aux767dux4e92ux52a8ux8868ohtdd4h-ux5206ux88c2ux53c2ux6570ux5bf9ux6bd4}

{\def\LTcaptype{none} % do not increment counter
\begin{longtable}[]{@{}
  >{\centering\arraybackslash}p{(\linewidth - 10\tabcolsep) * \real{0.1724}}
  >{\centering\arraybackslash}p{(\linewidth - 10\tabcolsep) * \real{0.1724}}
  >{\centering\arraybackslash}p{(\linewidth - 10\tabcolsep) * \real{0.1724}}
  >{\raggedright\arraybackslash}p{(\linewidth - 10\tabcolsep) * \real{0.1379}}
  >{\centering\arraybackslash}p{(\linewidth - 10\tabcolsep) * \real{0.1724}}
  >{\centering\arraybackslash}p{(\linewidth - 10\tabcolsep) * \real{0.1724}}@{}}
\toprule\noalign{}
\begin{minipage}[b]{\linewidth}\centering
配位场
\end{minipage} & \begin{minipage}[b]{\linewidth}\centering
配位数
\end{minipage} & \begin{minipage}[b]{\linewidth}\centering
高点群
\end{minipage} & \begin{minipage}[b]{\linewidth}\raggedright
分裂模式
\end{minipage} & \begin{minipage}[b]{\linewidth}\centering
Δ 表达式
\end{minipage} & \begin{minipage}[b]{\linewidth}\centering
Δ 典型值
\end{minipage} \\
\midrule\noalign{}
\endhead
\bottomrule\noalign{}
\endlastfoot
\textbf{八面体} Oh & 6 & \(O_h\) & \(t_{2g} \downarrow\)，\(e_g \uparrow\) & \(\Delta_o\) & 基准 \\
\textbf{四面体} Td & 4 & \(T_d\) & \_\_\_\_\_\_ & \$\Delta\emph{t = \$}\_\_\_\_\_ & ≈ \_\_\_\_\_\_ \\
\textbf{正方形} D4h & 4 & \(D_{4h}\) & \_\_\_\_\_\_ & \(\Delta_{sp} \approx\)\_\_\_\_\_\_ & ≈ \_\_\_\_\_\_ \\
\end{longtable}
}

（课堂填空：学生自行填写各场分裂模式和 Δ 表达式）

\paragraph{高自旋 vs 低自旋}\label{ux9ad8ux81eaux65cb-vs-ux4f4eux81eaux65cb}

对于 \(d^4 \sim d^7\) 构型，电子排布有两种可能：

\begin{itemize}
\tightlist
\item
  \textbf{高自旋}：\(\Delta < P\)（弱场配体）→ 电子优先分占各轨道（洪特规则）
\item
  \textbf{低自旋}：\(\Delta > P\)（强场配体）→ 电子优先填满低能级（配对）
\end{itemize}

\begin{quote}
🧠 \textbf{认知冲突}：「d¹-d3 和 d⁸-d¹⁰ 怎么不分高/低自旋？」------d¹-d3 电子数太少，还不足以塞满 t2g 轨道，没有''配对 vs 上 e\_g''的选择；d⁸-d¹⁰ 电子数太多，无论如何 e\_g 上必须有电子。\textbf{只有 d⁴-d⁷ 才有高低自旋的选择空间。}
\textbf{速记口诀}：「d⁴-d⁷ 才需选，d¹-d3 高自旋，d⁸-d¹⁰ 低自旋（八面体场）。」
\end{quote}

\textbf{光谱化学序列}（配体场强由弱到强）：

\[
I^- < Br^- < S^{2- < SCN^- < Cl^- < OH^- < C_2O_4^{2-} < H_2O < NH_3 < en < NO_2^- < CN^- < CO}
\]

\begin{quote}
🧠 \textbf{认知冲突}：「为什么 CO 是中性的，场强却比带负电的 I- 强得多？」------纯静电模型预言阴离子配体应当场强更大（负电吸引正电中心离子），但事实正好相反。原因是：CO 不仅有 σ 给电子（C 上的孤对电子给金属），还有 \textbf{π 反馈}（金属的 d 电子反馈到 CO 的 π* 空轨道）。这个 π 反馈增强了金属-配体键。I- 是 \textbf{π 给体}------它的 π 电子和金属的 d 电子排斥------减弱了 Δ。
\textbf{配体三分类}：
- \textbf{σ 给体}（NH3, en）：只给 σ 电子 → Δ 中等
- \textbf{π 给体}（I-, Br-, Cl-, F-）：有孤对 π 电子→推高 t2g→Δ 减小→\textbf{弱场}
- \textbf{π 受体}（CO, CN-, PPh3）：有空 π* 轨道→拉低 t2g→Δ 增大→\textbf{强场}
\end{quote}

\begin{quote}
🗣️ \textbf{课堂原话}：``光谱化学序列是整个配位化学的『元素周期表』------必须能背诵到 H2O 前后。先用弱场/强场两端记住，再向中间填：I- 最弱，CO 最强，NH3 在中间。''
\end{quote}

\begin{figure}[H]\centering
\includegraphics[width=0.9\textwidth,keepaspectratio]{spectrochemical_series.png}
\caption{图3c 光谱化学序列配体排序 — 从左至右场强递增：π 给体（I-, Br-, Cl-, F-, 红色）→ σ 给体（OH-, H2O, 黄色）→ σ 给体+螯合（NH3, en, 蓝色）→ π 受体（NO2-, CN-, CO, 绿色）。配体三分类决定 Δ 大小}
\end{figure}

{\def\LTcaptype{none} % do not increment counter
\begin{longtable}[]{@{}lll@{}}
\toprule\noalign{}
场强 & 配体 & 自旋状态 \\
\midrule\noalign{}
\endhead
\bottomrule\noalign{}
\endlastfoot
\textbf{弱场} & I-, Br-, Cl-, F-, H2O & 高自旋（外轨型） \\
\textbf{中强场} & NH3, en & 依金属离子而定 \\
\textbf{强场} & CN-, CO, NO2- & 低自旋（内轨型） \\
\end{longtable}
}

\begin{figure}[H]\centering
\includegraphics[width=0.9\textwidth,keepaspectratio]{spin_state_decision.png}
\caption{图4 高自旋与低自旋的 d 电子排布 — Δ 与 P 的相对大小决定电子排布方式（以 Fe2+/Fe3+ 为例）}
\end{figure}

\begin{figure}[H]\centering
\includegraphics[width=0.9\textwidth,keepaspectratio]{spin_state_decision.png}
\caption{图4b 高/低自旋判断决策树 — 从 dⁿ 组态出发，经几何构型判断，再到 Δ 与 P 的比较，最终决定 HS/LS 状态。附磁矩 μ = √[n(n+2)] 速查}
\end{figure}

\paragraph{晶体场稳定化能（CFSE）}\label{ux6676ux4f53ux573aux7a33ux5b9aux5316ux80fdcfse}

CFSE 是 d 电子在配体场中排布后，相对于未分裂时的球形场，体系能量降低的值（越负越稳定）。

\[
CFSE_{Oh} = n_{t_{2g}} \times (-4\ Dq) + n_{e_g} \times (+6\ Dq)
\]

其中 \(n_{t_{2g}}\) 为 t2g 轨道上的电子数，\(n_{e_g}\) 为 e\_g 轨道上的电子数。

考虑成对能修正（低自旋时需要额外克服配对能）：

\[
CFSE_{Oh} = [n_{t_{2g}} \times (-4\ Dq) + n_{e_g} \times (+6\ Dq)] - mP
\]

其中 m 为\textbf{成对数增加值}（相对于自由离子），P 为成对能。

\begin{quote}
🧠 \textbf{认知冲突}：「CFSE 越负配合物越稳定，d⁵ 高自旋的 CFSE = 0------是不是它『没有』稳定化能？」------是的，但注意 CFSE 只是配合物稳定性的一个因素。螯合效应、溶剂化能、熵效应同样重要。CFSE 为零不代表配合物不稳定------比如 \([Mn(H_2O)_6]^{2+}\)（d⁵ 高自旋，CFSE = 0）在溶液中相当稳定。
\end{quote}

\begin{figure}[H]\centering
\includegraphics[width=0.9\textwidth,keepaspectratio]{cfse_calculation.png}
\caption{图4c CFSE 计算四步流程 — Step 1 dⁿ→Step 2 几何构型（Oh/Td/D4h）→Step 3 按分裂能级填充电子→Step 4 代入 CFSE 公式。右侧附高低自旋判断要点}
\end{figure}

\paragraph{d¹-d¹⁰ CFSE 计算表（八面体场）}\label{duxb9-duxb9ux2070-cfse-ux8ba1ux7b97ux8868ux516bux9762ux4f53ux573a}

{\def\LTcaptype{none} % do not increment counter
\begin{longtable}[]{@{}
  >{\centering\arraybackslash}p{(\linewidth - 10\tabcolsep) * \real{0.1852}}
  >{\raggedright\arraybackslash}p{(\linewidth - 10\tabcolsep) * \real{0.1481}}
  >{\centering\arraybackslash}p{(\linewidth - 10\tabcolsep) * \real{0.1852}}
  >{\raggedright\arraybackslash}p{(\linewidth - 10\tabcolsep) * \real{0.1481}}
  >{\centering\arraybackslash}p{(\linewidth - 10\tabcolsep) * \real{0.1852}}
  >{\raggedright\arraybackslash}p{(\linewidth - 10\tabcolsep) * \real{0.1481}}@{}}
\toprule\noalign{}
\begin{minipage}[b]{\linewidth}\centering
dⁿ
\end{minipage} & \begin{minipage}[b]{\linewidth}\raggedright
高自旋排布
\end{minipage} & \begin{minipage}[b]{\linewidth}\centering
HS-CFSE (Dq)
\end{minipage} & \begin{minipage}[b]{\linewidth}\raggedright
低自旋排布
\end{minipage} & \begin{minipage}[b]{\linewidth}\centering
LS-CFSE (Dq)
\end{minipage} & \begin{minipage}[b]{\linewidth}\raggedright
说明
\end{minipage} \\
\midrule\noalign{}
\endhead
\bottomrule\noalign{}
\endlastfoot
d⁰ & --- & 0 & --- & 0 & 无 d 电子 \\
d¹ & \(t_{2g}^1\) & -4 & --- & --- & 仅一种 \\
d2 & \(t_{2g}^2\) & -8 & --- & --- & 仅一种 \\
d3 & \(t_{2g}^3\) & -12 & --- & --- & 仅一种 \\
d⁴ & \(t_{2g}^3 e_g^1\) & -6 & \(t_{2g}^4\) & -16 + P & 高低自旋分水岭 \\
d⁵ & \(t_{2g}^3 e_g^2\) & \textbf{0} & \(t_{2g}^5\) & -20 + 2P & HS-CFSE=0 的特殊情况 \\
d⁶ & \(t_{2g}^4 e_g^2\) & -4 & \(t_{2g}^6\) & \textbf{-24 + 2P} & LS 稳定化能最大 \\
d⁷ & \(t_{2g}^5 e_g^2\) & -8 & \(t_{2g}^6 e_g^1\) & -18 + P & 最后的分水岭 \\
d⁸ & \(t_{2g}^6 e_g^2\) & -12 & --- & --- & 仅一种 \\
d⁹ & \(t_{2g}^6 e_g^3\) & -6 & --- & --- & 仅一种（有 JT 畸变） \\
d¹⁰ & \(t_{2g}^6 e_g^4\) & 0 & --- & --- & 全满 \\
\end{longtable}
}

\begin{quote}
⚠️ \textbf{关键发现}：
- d⁶ 低自旋的 CFSE = -24 Dq------这是所有 dⁿ 中\textbf{绝对值最大的}，解释了为什么 Co3+（d⁶）的配合物特别稳定
- d⁵ 高自旋 CFSE = 0------高自旋和低自旋唯一相同的组态（\(t_{2g}^3\) 的 -12 刚好被 \(e_g^2\) 的 +12 抵消）
- d⁴-d⁷ 是唯一有高低自旋选择的区间------记住这个范围！
- 四面体场中，因分裂能小（\(\Delta_t \approx 4/9 \Delta_o\)），通常只有\textbf{高自旋}
\end{quote}

\paragraph{CFSE 的应用}\label{cfse-ux7684ux5e94ux7528}

\textbf{① 稳定性比较}：同构型配合物中，CFSE 越负 → 配合物越稳定。

\textbf{② 解释 Jahn-Teller 效应}：\(d^9\)（如 Cu2+）、\(d^7\)（高自旋）、\(d^4\)（高自旋）构型发生畸变------基态简并的八面体配合物通过几何畸变消除简并，进一步降低能量。

\textbf{Cu2+ 配合物为什么几乎总是拉长的八面体？}------d⁹ 的 e\_g 轨道上有 3 个电子（\(e_g^3\)），相当于 1 个空穴。拉长 z 轴配体后，d\_z2 能量降低，空穴留在 dₓ2₋ᵧ2 上 → 额外稳定化能。

\textbf{③ 解释配合物的颜色}：d-d 跃迁------电子从 t2g 激发到 e\_g，吸收特定波长的可见光。吸收波长 λ 与 Δ 的关系：

\[
\Delta_o = h\nu = \frac{hc}{\lambda}
\]

\begin{figure}[H]\centering
\includegraphics[width=0.9\textwidth,keepaspectratio]{d-d-transition.jpg}
\caption{图5 d-d 跃迁示意图 — 电子从 t2g 激发到 e_g，吸收特定波长可见光}
\end{figure}

\begin{figure}[H]\centering
\includegraphics[width=0.9\textwidth,keepaspectratio]{14-7-ti-h2o-6-d-d-transition-absorption-spectrum.jpg}
\caption{图6 [Ti(H2O)6]3+ 的 d-d 跃迁吸收光谱 — d¹ 体系在约 490 nm 处有一个吸收峰，配合物呈紫红色（普通化学原理第4版 图14.7）}
\end{figure}

\begin{quote}
📝 \textbf{练一练}：已知 {[}Ti(H2O)6{]}3+ 的 \(\lambda_{\max} = 490\ nm\)。
(1) 计算 \(\Delta_o\)（单位 kJ/mol）
(2) 预测配合物的颜色
\end{quote}

\textbf{参考答案}：
(1) \(\Delta_o = N_A \times hc / \lambda = 6.02\times10^{23} \times 6.626\times10^{-34} \times 3.0\times10^8 / (490\times10^{-9}) \approx 239\ kJ/mol\)
(2) 吸收绿光（\textasciitilde490 nm），呈现\textbf{紫红色}（补色）

相关 KP：晶体场理论 · 高自旋与低自旋 · 光谱化学序列 · d轨道分裂 · Jahn-Teller效应 · 磁矩

\subsubsection{2.4 配合物的异构现象}\label{ux914dux5408ux7269ux7684ux5f02ux6784ux73b0ux8c61}

配合物的异构现象比有机化学更丰富------不仅包括结构异构（化学式相同但成键方式不同），还包括几何异构（空间排列不同）和光学异构（手性）。

\paragraph{空白互动表：异构体分类对比}\label{ux7a7aux767dux4e92ux52a8ux8868ux5f02ux6784ux4f53ux5206ux7c7bux5bf9ux6bd4}

{\def\LTcaptype{none} % do not increment counter
\begin{longtable}[]{@{}
  >{\centering\arraybackslash}p{(\linewidth - 6\tabcolsep) * \real{0.2778}}
  >{\raggedright\arraybackslash}p{(\linewidth - 6\tabcolsep) * \real{0.2222}}
  >{\raggedright\arraybackslash}p{(\linewidth - 6\tabcolsep) * \real{0.2222}}
  >{\centering\arraybackslash}p{(\linewidth - 6\tabcolsep) * \real{0.2778}}@{}}
\toprule\noalign{}
\begin{minipage}[b]{\linewidth}\centering
异构类型
\end{minipage} & \begin{minipage}[b]{\linewidth}\raggedright
识别要点
\end{minipage} & \begin{minipage}[b]{\linewidth}\raggedright
经典实例
\end{minipage} & \begin{minipage}[b]{\linewidth}\centering
手性可能
\end{minipage} \\
\midrule\noalign{}
\endhead
\bottomrule\noalign{}
\endlastfoot
\textbf{结构异构} & 原子连接方式不同 & \_\_\_\_\_\_ & 无关 \\
├ 电离异构 & 内外界成分交换 & \_\_\_\_\_\_ & --- \\
├ 水合异构 & H2O 在内界/外界分布不同 & \_\_\_\_\_\_ & --- \\
├ 键合异构 & 两可配体以不同原子配位 & \_\_\_\_\_\_ & --- \\
└ 配位异构 & 配正/负离子之间交换配体 & \_\_\_\_\_\_ & --- \\
\textbf{几何异构} & 配体空间排列不同（cis/trans, fac/mer） & \_\_\_\_\_\_ & 顺式可能有 \\
\textbf{光学异构} & 分子有手性（无σ/无i/无S4）→对映体 & \_\_\_\_\_\_ & 必然有 \\
\end{longtable}
}

（课堂填空：学生填写各类型的实例）

\begin{figure}[H]\centering
\includegraphics[width=0.9\textwidth,keepaspectratio]{isomer_classification.png}
\caption{图5 配合物异构现象分类体系 — 三类异构体：结构异构（电离/水合/键合/配位）、几何异构（cis/trans、fac/mer）、光学异构（Δ/Λ 对映体）}
\end{figure}

\paragraph{（1）结构异构（构造异构）}\label{ux7ed3ux6784ux5f02ux6784ux6784ux9020ux5f02ux6784}

{\def\LTcaptype{none} % do not increment counter
\begin{longtable}[]{@{}
  >{\raggedright\arraybackslash}p{(\linewidth - 4\tabcolsep) * \real{0.3333}}
  >{\raggedright\arraybackslash}p{(\linewidth - 4\tabcolsep) * \real{0.3333}}
  >{\raggedright\arraybackslash}p{(\linewidth - 4\tabcolsep) * \real{0.3333}}@{}}
\toprule\noalign{}
\begin{minipage}[b]{\linewidth}\raggedright
类型
\end{minipage} & \begin{minipage}[b]{\linewidth}\raggedright
特征
\end{minipage} & \begin{minipage}[b]{\linewidth}\raggedright
示例
\end{minipage} \\
\midrule\noalign{}
\endhead
\bottomrule\noalign{}
\endlastfoot
\textbf{电离异构} & 内外界交换成分，电离出不同离子 & \([CoBr(NH_3)_5]SO_4\)（SO42-在外界，加 Ba2+ 有沉淀）vs \([CoSO_4(NH_3)_5]Br\)（Br-在外界，加 Ag+ 有沉淀） \\
\textbf{水合异构} & 水分子在内界和外界的分布不同 & \([Cr(H_2O)_6]Cl_3\)（紫色）vs \([CrCl(H_2O)_5]Cl_2\cdot H_2O\)（蓝绿色）------颜色不同！ \\
\textbf{键合异构} & 两可配体通过不同配位原子连接 & \([Co(NO_2)(NH_3)_5]Cl_2\)（硝基，N-配位）vs \([Co(ONO)(NH_3)_5]Cl_2\)（亚硝酸根，O-配位） \\
\textbf{配位异构} & 配正/负离子之间交换配体 & \([Co(NH_3)_6][Cr(CN)_6]\) vs \([Cr(NH_3)_6][Co(CN)_6]\) \\
\end{longtable}
}

\paragraph{（2）几何异构（顺反异构 + fac/mer）}\label{ux51e0ux4f55ux5f02ux6784ux987aux53cdux5f02ux6784-facmer}

出现在配位数 ≥ 4 的配合物中，最常见于\textbf{平面正方形}和\textbf{八面体}配合物。

\textbf{平面正方形 \(Ma_2c_2\) 型}（如 \([Pt(NH_3)_2Cl_2]\)）：
- \textbf{顺式}（cis）：相同配体相邻 → 极性分子 → 有抗癌活性（顺铂）
- \textbf{反式}（trans）：相同配体相对 → 非极性分子 → 无抗癌活性

\textbf{八面体 \(Ma_4e_2\) 型}（如 \([Co(NH_3)_4Cl_2]^+\)）：
- 顺式（cis/1,2-）：2 个 Cl 相邻 → 有光学活性
- 反式（trans/1,6-）：2 个 Cl 相对 → 有对称面 → 无光学活性

\textbf{八面体 \(Ma_3b_3\) 型}（如 \([CoCl_3(NH_3)_3]\)）：
- \textbf{面式（fac）}：3 个相同配体共占八面体的一个面 → \textbf{有手性}（一对对映体）
- \textbf{经式（mer）}：3 个相同配体共占八面体的一条经线 → 无手性

\begin{quote}
🧠 \textbf{认知冲突}：「Ma3b3 型八面体配合物有几种异构体？」------大多数学生只画出 fac 和 mer 两种几何异构就停了。但 fac 结构无对称面→有手性→存在\textbf{一对对映体}！所以总共是 3 种：反式 1 种 + 顺式一对（2 种）。\textbf{枚举异构体后别忘了检查手性！}
\end{quote}

\begin{quote}
🗣️ \textbf{课堂原话}：「先定骨架，再排配体，最后查手性------这是枚举异构体的『固定-变换法』。配体种类越多，八面体的异构体越复杂，Ma2b2c2 有 6 种异构体（含光学异构）。」
\end{quote}

\paragraph{（3）光学异构（镜面异构/手性异构）}\label{ux5149ux5b66ux5f02ux6784ux955cux9762ux5f02ux6784ux624bux6027ux5f02ux6784}

\textbf{条件}：分子无对称面、无对称中心、无交替轴（S4）→ 手性分子 → 存在一对对映体。

\textbf{口诀}：「无 σ 无 i 无 S4 则 Chiral」

\textbf{常见手性配合物}：
- \([M(AA)_3]\) 型（AA 为双齿配体）：如 \([Co(en)_3]^{3+}\)、\([Cr(ox)_3]^{3-}\)------有 Λ 和 Δ 两种对映体
- \(cis-[M(AA)_2X_2]\) 型：如 \(cis-[Co(en)_2Cl_2]^+\)
- \textbf{四面体 \(Mabcd\)} 型（四个不同配体）

\begin{quote}
💡 \textbf{判断技巧}：若配合物存在对称面或对称中心 → 无光学异构。\(trans-[M(AA)_2X_2]\) 有对称面 → 无光学异构；\(cis-\) 无对称面 → 有光学异构。
\end{quote}

相关 KP：配合物 · 配位异构 · 手性

\subsubsection{2.5 常见误区与认知陷阱}\label{ux5e38ux89c1ux8befux533aux4e0eux8ba4ux77e5ux9677ux9631}

\begin{quote}
⚠️ \textbf{误区 1}：「配合物一定有外界。」
\textbf{纠正}：\([Ni(CO)_4]\)、\([Fe(CO)_5]\) 等中性配合物只有内界，没有外界。配合物整体呈电中性时不需要外界离子。
\end{quote}

\begin{quote}
⚠️ \textbf{误区 2}：「内轨型配合物一定有磁性变化。」
\textbf{纠正}：d¹\textasciitilde d3 构型本身有空的 d 轨道，形成内轨型配合物时 d 电子排布不变，磁性不变。例如 \([Cr(NH_3)_6]^{3+}\)（Cr3+ d3）为内轨型（d2sp3），但仍有 3 个未成对电子。
\end{quote}

\begin{quote}
⚠️ \textbf{误区 3}：「高自旋 = 外轨型，低自旋 = 内轨型，两者完全等同。」
\end{quote}

\begin{quote}
⚠️ \textbf{误区 4}：「所有八面体配合物都有高/低自旋之分。」
\textbf{纠正}：d¹\textasciitilde d3 只有一种排布（电子数太少，无法选择）；d⁸\textasciitilde d¹⁰ 也只有一种排布。\textbf{只有 d⁴\textasciitilde d⁷ 才有两种可能。}
\end{quote}

\begin{quote}
⚠️ \textbf{误区 5}：「命名时配体顺序可以随意。」
\textbf{纠正}：IUPAC 规则严格规定：先负离子→后中性分子→再按配位原子字母顺序。这个是命名题主要扣分点。
\end{quote}

\begin{quote}
⚠️ \textbf{误区 6}：「CFSE 越大配合物就一定越稳定。」
\textbf{纠正}：CFSE 是重要因素，但还需考虑螯合效应（熵增）、溶剂化能、配体间排斥等因素。例如 \([Ni(en)_3]^{2+}\) 的 CFSE 比 \([Ni(NH_3)_6]^{2+}\) 大不了很多，但前者\textbf{稳定得多}------这是螯合效应（熵增）的贡献。
\end{quote}

\subsection{三、典型例题}\label{ux4e09ux5178ux578bux4f8bux9898}

\subsubsection{例 1 {[}配合物命名与组成{]} ⭐⭐}\label{ux4f8b-1-ux914dux5408ux7269ux547dux540dux4e0eux7ec4ux6210}

\textbf{题目}：写出下列配合物的名称，并指出中心离子、配体、配位数：
(1) \(K_3[Fe(CN)_6]\) (2) \([Co(NH_3)_5(H_2O)]Cl_3\) (3) \(H_2[SiF_6]\)

\textbf{完整解答}：

{\def\LTcaptype{none} % do not increment counter
\begin{longtable}[]{@{}
  >{\raggedright\arraybackslash}p{(\linewidth - 8\tabcolsep) * \real{0.1667}}
  >{\centering\arraybackslash}p{(\linewidth - 8\tabcolsep) * \real{0.2083}}
  >{\centering\arraybackslash}p{(\linewidth - 8\tabcolsep) * \real{0.2083}}
  >{\centering\arraybackslash}p{(\linewidth - 8\tabcolsep) * \real{0.2083}}
  >{\centering\arraybackslash}p{(\linewidth - 8\tabcolsep) * \real{0.2083}}@{}}
\toprule\noalign{}
\begin{minipage}[b]{\linewidth}\raggedright
配合物
\end{minipage} & \begin{minipage}[b]{\linewidth}\centering
名称
\end{minipage} & \begin{minipage}[b]{\linewidth}\centering
中心离子
\end{minipage} & \begin{minipage}[b]{\linewidth}\centering
配体
\end{minipage} & \begin{minipage}[b]{\linewidth}\centering
配位数
\end{minipage} \\
\midrule\noalign{}
\endhead
\bottomrule\noalign{}
\endlastfoot
\(K_3[Fe(CN)_6]\) & 六氰合铁(Ⅲ)酸钾 & Fe3+ & 6 × CN- & 6 \\
\([Co(NH_3)_5(H_2O)]Cl_3\) & 三氯化五氨·一水合钴(Ⅲ) & Co3+ & 5 × NH3 + 1 × H2O & 6 \\
\(H_2[SiF_6]\) & 六氟合硅(Ⅳ)酸 & Si⁴+ & 6 × F- & 6 \\
\end{longtable}
}

\subsubsection{例 2 {[}内/外轨型+磁性判断{]} ⭐⭐⭐}\label{ux4f8b-2-ux5185ux5916ux8f68ux578bux78c1ux6027ux5224ux65ad}

\textbf{题目}：实验测得 \([CoF_6]^{3-}\) 的磁矩为 5.26 B.M.，\([Co(CN)_6]^{3-}\) 的磁矩为 0 B.M.。用价键理论推测两者的杂化类型、空间构型和内/外轨类型。

\textbf{思路分析}：由磁矩 → 单电子数 n → d 电子排布是否改变 → 杂化类型

\textbf{完整解答}：

\textbf{\([CoF_6]^{3-}\)}：
- \(\mu = \sqrt{n(n+2)} = 5.26 \rightarrow n = 4\)
- Co3+ 为 \(d^6\)，自由离子有 4 个未成对电子
- 形成配合物后 n 仍为 4 → d 电子排布\textbf{未改变}
- 使用外层 4d 轨道 → \textbf{sp3d2 杂化} → 正八面体 → \textbf{外轨型} → 高自旋

\textbf{\([Co(CN)_6]^{3-}\)}：
- \(\mu = 0 \rightarrow n = 0\)
- 形成配合物后 n 从 4 变为 0 → d 电子\textbf{重排归并}
- 空出内层 3d 轨道 → \textbf{d2sp3 杂化} → 正八面体 → \textbf{内轨型} → 低自旋

\subsubsection{例 3 {[}EAN 18 电子规则{]} ⭐⭐⭐}\label{ux4f8b-3-ean-18-ux7535ux5b50ux89c4ux5219}

\textbf{题目}：判断下列配合物是否符合 18 电子规则：\([Ni(CO)_4]\)、\([Mn(CO)_6]^+\)、\([Fe_2(CO)_9]\)

\textbf{思路分析}：中心原子价电子数 + 配体提供的电子数 = 18？

\textbf{完整解答}：

{\def\LTcaptype{none} % do not increment counter
\begin{longtable}[]{@{}
  >{\raggedright\arraybackslash}p{(\linewidth - 8\tabcolsep) * \real{0.1667}}
  >{\centering\arraybackslash}p{(\linewidth - 8\tabcolsep) * \real{0.2083}}
  >{\centering\arraybackslash}p{(\linewidth - 8\tabcolsep) * \real{0.2083}}
  >{\centering\arraybackslash}p{(\linewidth - 8\tabcolsep) * \real{0.2083}}
  >{\centering\arraybackslash}p{(\linewidth - 8\tabcolsep) * \real{0.2083}}@{}}
\toprule\noalign{}
\begin{minipage}[b]{\linewidth}\raggedright
配合物
\end{minipage} & \begin{minipage}[b]{\linewidth}\centering
中心原子价电子
\end{minipage} & \begin{minipage}[b]{\linewidth}\centering
配体贡献
\end{minipage} & \begin{minipage}[b]{\linewidth}\centering
总和
\end{minipage} & \begin{minipage}[b]{\linewidth}\centering
符合？
\end{minipage} \\
\midrule\noalign{}
\endhead
\bottomrule\noalign{}
\endlastfoot
\([Ni(CO)_4]\) & Ni(0): 10e- & 4 × CO = 8e- & 18 & ✅ \\
\([Mn(CO)_6]^+\) & Mn+: 6e- & 6 × CO = 12e- & 18 & ✅ \\
\([Fe_2(CO)_9]\) & Fe(0): 8e- & 3×桥CO=3e- + 3×端CO=6e- + 1×Fe-Fe=1e- & 18 & ✅ \\
\end{longtable}
}

\textbf{关键}：\([Fe_2(CO)_9]\) 中有 3 个桥联 CO（每个为每个 Fe 贡献 1e-），3 个端基 CO（每个 2e-），以及 1 条 Fe---Fe 键（1e-），合计 \(8 + 3 + 6 + 1 = 18\)。

\subsubsection{例 4 {[}CFSE 计算 + 高/低自旋判断{]} ⭐⭐⭐⭐}\label{ux4f8b-4-cfse-ux8ba1ux7b97-ux9ad8ux4f4eux81eaux65cbux5224ux65ad}

\textbf{题目}：比较 \([Fe(H_2O)_6]^{2+}\) 和 \([Fe(CN)_6]^{4-}\) 的 CFSE 和自旋状态。已知 Fe2+ 的 P = 179.4 kJ/mol，\(\Delta_o(H_2O) = 124.4\ kJ/mol\)，\(\Delta_o(CN^-) = 394.7\ kJ/mol\)。

\textbf{思路分析}：比较 Δ 和 P → 判断高/低自旋 → 算 CFSE

\textbf{完整解答}：

\textbf{\([Fe(H_2O)_6]^{2+}\)}：Δ = 124.4 \textless{} P = 179.4 → \textbf{高自旋}
- d⁶ 高自旋：\(t_{2g^4 e_g^2}\)（4 个电子在 t2g，2 个在 e\_g）
- CFSE = \(4(-4\ Dq) + 2(+6\ Dq) = -4\ Dq\)

\textbf{\([Fe(CN)_6]^{4-}\)}：Δ = 394.7 \textgreater{} P = 179.4 → \textbf{低自旋}
- d⁶ 低自旋：\(t_{2g^6 e_g^0}\)（6 个电子全部在 t2g）
- CFSE = \(6(-4\ Dq) + 0(+6\ Dq) = -24\ Dq\)

\textbf{结论}：\([Fe(CN)_6]^{4-}\) 的 CFSE 远大于 \([Fe(H_2O)_6]^{2+}\)（差值 20 Dq），因此前者更稳定。配合物颜色的差异也源于 Δ 的不同------前者吸收蓝紫色光（Δ 大→\(\lambda\) 短），后者吸收红光（Δ 小→\(\lambda\) 长）。

\subsubsection{例 5 {[}异构体判断与枚举{]} ⭐⭐⭐⭐}\label{ux4f8b-5-ux5f02ux6784ux4f53ux5224ux65adux4e0eux679aux4e3e}

\textbf{题目}：画出 \([Co(en)_2Cl_2]^+\) 的所有异构体，并判断哪些有光学活性。

\textbf{思路分析}：en 是双齿配体，两个 Cl 的位置决定顺/反 → 顺式有光学异构

\textbf{完整解答}：

\textbf{反式异构体}（trans）：
- 2 个 Cl 处于对位
- 分子有对称面 → \textbf{无光学活性}

\textbf{顺式异构体}（cis）：
- 2 个 Cl 处于邻位
- 分子无对称面、无对称中心 → \textbf{存在一对对映体}（左旋 + 右旋）

\textbf{总计}：1 个反式异构体 + 1 对顺式光学异构体 = \textbf{3 种异构体}

\subsubsection{例 6 {[}磁性 + 颜色 + 构型综合判断{]} ⭐⭐⭐⭐}\label{ux4f8b-6-ux78c1ux6027-ux989cux8272-ux6784ux578bux7efcux5408ux5224ux65ad}

\textbf{题目}：某含 Fe 配合物的化学式为 \(K_3[FeF_6]\)，实验测得 \(\mu = 5.92\ B.M.\)。
(1) 判断 Fe 的氧化态和 dⁿ 组态；
(2) 判断高/低自旋并写出 d 电子排布；
(3) 计算 CFSE（以 Dq 为单位）；
(4) 预测该配合物的颜色。

\textbf{完整解答}：

\begin{enumerate}
\def\labelenumi{(\arabic{enumi})}
\item
  K3 → 3 个 K+ → 配阴离子为 {[}FeF6{]}3- → Fe 氧化态为 +3 → Fe3+ 为 d⁵
\item
  \(\mu = \sqrt{n(n+2)} = 5.92 \rightarrow n = 5\) → 5 个未成对电子 → \textbf{高自旋}
\end{enumerate}

\begin{itemize}
\tightlist
\item
  排布：\(t_{2g^3 e_g^2}\)
\end{itemize}

\begin{enumerate}
\def\labelenumi{(\arabic{enumi})}
\setcounter{enumi}{2}
\tightlist
\item
  CFSE = \(3(-4\ Dq) + 2(+6\ Dq) = 0\ Dq\)
\end{enumerate}

\begin{itemize}
\tightlist
\item
  d⁵ 高自旋的 CFSE 恰好为零（\(t_{2g}^3\) 的 -12 Dq 与 \(e_g^2\) 的 +12 Dq 抵消）
\end{itemize}

\begin{enumerate}
\def\labelenumi{(\arabic{enumi})}
\setcounter{enumi}{3}
\tightlist
\item
  F- 是弱场配体 → Δ 小 → 吸收红光（\textasciitilde700 nm）→ 配合物呈\textbf{浅绿色}（补色）
  实际上 K3{[}FeF6{]} 为浅绿色晶体，与预测一致。
\end{enumerate}

\subsection{四、拓展与联系}\label{ux56dbux62d3ux5c55ux4e0eux8054ux7cfb}

\subsubsection{4.1 与竞赛真题的对接}\label{ux4e0eux7adeux8d5bux771fux9898ux7684ux5bf9ux63a5}

{\def\LTcaptype{none} % do not increment counter
\begin{longtable}[]{@{}
  >{\raggedright\arraybackslash}p{(\linewidth - 4\tabcolsep) * \real{0.3077}}
  >{\centering\arraybackslash}p{(\linewidth - 4\tabcolsep) * \real{0.3846}}
  >{\raggedright\arraybackslash}p{(\linewidth - 4\tabcolsep) * \real{0.3077}}@{}}
\toprule\noalign{}
\begin{minipage}[b]{\linewidth}\raggedright
考点
\end{minipage} & \begin{minipage}[b]{\linewidth}\centering
频率
\end{minipage} & \begin{minipage}[b]{\linewidth}\raggedright
说明
\end{minipage} \\
\midrule\noalign{}
\endhead
\bottomrule\noalign{}
\endlastfoot
价键理论 + 内/外轨 + 磁性 & 每年 1-2 题 & 常给出磁矩数据反推结构 \\
CFSE 计算 + 高/低自旋 & 每年 1 题 & 常与颜色、稳定性结合考查 \\
几何异构/光学异构 & 隔年 1 题 & 常以 \([M(AA)_2X_2]\) 或 \([M(AA)_3]\) 形式出现 \\
EAN 18 电子规则 & 隔年 1 题 & 金属有机配合物推断 \\
配合物命名 & 每年 1 题 & 基础分，不可失 \\
\end{longtable}
}

\subsubsection{4.2 跨专题联系}\label{ux8de8ux4e13ux9898ux8054ux7cfb}

\begin{itemize}
\tightlist
\item
  杂化轨道理论 → 价键理论的基础
\item
  分子轨道理论 → 配位场理论的更完整处理（ML6 σ + π 成键）
\item
  晶体场理论 → Jahn-Teller 效应 → 配合物颜色与光谱
\item
  18电子规则 → 金属有机化学 → 羰基配合物
\item
  配位异构 → 手性 → 旋光性
\item
  过渡金属通性 → 配合物的元素化学特性
\end{itemize}

\subsubsection{4.3 深化方向（后续轮次展开）}\label{ux6df1ux5316ux65b9ux5411ux540eux7eedux8f6eux6b21ux5c55ux5f00}

🔷 本节建立配位化学的入门基础，以下内容将在后续轮次深化：
- \textbf{分子轨道理论处理配合物}（AOM、σ/π 成键、光谱化学序列的理论解释）------第三轮
- \textbf{Jahn-Teller 效应的定量分析}（畸变能与结构预测）------第三轮结构深化
- \textbf{配位催化循环}（Wacker 氧化、氢甲酰化）------第三轮金属有机
- \textbf{生物无机配合物}（卟啉、铁硫簇、固氮酶）------决赛拓展

\subsection{五、方法总结}\label{ux4e94ux65b9ux6cd5ux603bux7ed3}

\subsubsection{5.1 配合物分析四步法}\label{ux914dux5408ux7269ux5206ux6790ux56dbux6b65ux6cd5}

\begin{Shaded}
\begin{Highlighting}[]
\NormalTok{给定配合物信息（化学式、磁性、颜色等）}
\NormalTok{    ↓}
\NormalTok{Step 1: 识别}
\NormalTok{   确定中心离子氧化态 + dⁿ 组态 + 配体种类和数目 → 配位数}
\NormalTok{    ☐ 氧化态计算正确，注意配合物总电荷}
\NormalTok{    ☐ 配体 denticity 正确（单齿/双齿/多齿）}
\NormalTok{    ↓}
\NormalTok{Step 2: 构型}
\NormalTok{   由配位数和 dⁿ 判断构型（八面体/四面体/平面正方形）}
\NormalTok{   考虑 Jahn{-}Teller 畸变可能（d⁴ HS / d⁷ LS / d⁹）}
\NormalTok{    ☐ d⁰ 无 CFSE 偏好，d⁸+强场→平面正方形}
\NormalTok{    ↓}
\NormalTok{Step 3: 电子排布}
\NormalTok{   确定高/低自旋（看配体场强 vs dⁿ 范围）}
\NormalTok{   写 t2gˣ e\_gʸ → 算 CFSE → 算未成对电子数 n → 磁矩 μ = √[n(n+2)]}
\NormalTok{    ☐ d⁴{-}d⁷ 才有高低自旋之分}
\NormalTok{    ☐ 4d/5d 几乎总是低自旋}
\NormalTok{    ☐ CFSE = n(t2g)({-}4) + n(e\_g)(+6) Dq（八面体场）}
\NormalTok{    ↓}
\NormalTok{Step 4: 性质推断}
\NormalTok{   颜色（d{-}d 跃迁吸收补色）→ 稳定性 → 异构体计数}
\NormalTok{    ☐ d⁰/d¹⁰ 配合物通常无色}
\NormalTok{    ☐ 枚举异构体后检查手性}
\end{Highlighting}
\end{Shaded}

\subsubsection{5.2 核心速查表}\label{ux6838ux5fc3ux901fux67e5ux8868}

{\def\LTcaptype{none} % do not increment counter
\begin{longtable}[]{@{}
  >{\raggedright\arraybackslash}p{(\linewidth - 4\tabcolsep) * \real{0.3333}}
  >{\raggedright\arraybackslash}p{(\linewidth - 4\tabcolsep) * \real{0.3333}}
  >{\raggedright\arraybackslash}p{(\linewidth - 4\tabcolsep) * \real{0.3333}}@{}}
\toprule\noalign{}
\begin{minipage}[b]{\linewidth}\raggedright
概念
\end{minipage} & \begin{minipage}[b]{\linewidth}\raggedright
关键公式/规律
\end{minipage} & \begin{minipage}[b]{\linewidth}\raggedright
注意点
\end{minipage} \\
\midrule\noalign{}
\endhead
\bottomrule\noalign{}
\endlastfoot
\textbf{配合物组成} & 内界 + 外界；中心离子 + 配体 & 中性配合物无外界 \\
\textbf{命名} & 配体数-配体名-``合''-中心离子-(氧化数) & 配体顺序：负→中，无→有，字母序 \\
\textbf{价键理论} & 杂化类型 → 空间构型 & 相同配位数可不同杂化→不同构型 \\
\textbf{内轨/外轨} & 内轨：(n-1)d 参与，d 电子重排；外轨：nd 参与，不重排 & 须磁矩实验判断 \\
\textbf{EAN 规则} & 金属价电子 + 配体电子 = 18 & NO 按 3e- 计；M---M 键 1e-/条 \\
\textbf{晶体场理论} & Δ \textgreater{} P → 低自旋；Δ \textless{} P → 高自旋 & d¹-3、d⁸-¹⁰ 无高低自旋之分 \\
\textbf{CFSE（Oh）} & \(n_{t_{2g}}(-4) + n_{e_g}(+6)\) Dq & 零场 (d⁰/d⁵高/d¹⁰弱场) → CFSE = 0 \\
\textbf{磁矩公式} & \(\mu = \sqrt{n(n+2)}\) B.M. & μ = 0 → 反磁性；μ \textgreater{} 0 → 顺磁性 \\
\textbf{光谱化学序列} & I- \textless{} Br- \textless{} Cl- \textless{} F- \textless{} H2O \textless{} NH3 \textless{} en \textless{} NO2- \textless{} CN- \textless{} CO & 记住两端：I- 最弱，CO 最强 \\
\textbf{异构体} & 结构异构 + 几何异构 + 光学异构 & 反式有对称面→无光学活性；fac 有手性 \\
\end{longtable}
}

\subsubsection{5.3 高/低自旋判断流程图}\label{ux9ad8ux4f4eux81eaux65cbux5224ux65adux6d41ux7a0bux56fe}

\begin{Shaded}
\begin{Highlighting}[]
\NormalTok{给定中心离子 dⁿ 组态 + 配体}
\NormalTok{    ↓}
\NormalTok{dⁿ 在 d⁴\textasciitilde{}d⁷ 范围内？}
\NormalTok{├─ 否（d⁰{-}d3 / d⁸{-}d¹⁰）→ 只有一种排布方式，无 HS/LS 之分}
\NormalTok{└─ 是 → 查配体在光谱化学序列中的位置}
\NormalTok{    ├─ 弱场配体（I{-}, Br{-}, Cl{-}, F{-}, H2O）→ Δ \textless{} P → 高自旋}
\NormalTok{    ├─ 中强场（NH3, en）→ 依金属而定}
\NormalTok{    └─ 强场配体（CN{-}, CO, NO2{-}）→ Δ \textgreater{} P → 低自旋}
\NormalTok{        ↓}
\NormalTok{写 d 电子排布 → 算 CFSE → 算 μ}
\end{Highlighting}
\end{Shaded}

\subsection{六、练习题}\label{ux516dux7ec3ux4e60ux9898}

\subsubsection{6.1 基础巩固（⭐）}\label{ux57faux7840ux5de9ux56fa}

\begin{enumerate}
\def\labelenumi{\arabic{enumi}.}
\item
  命名下列配合物，并指出中心离子的氧化数和配位数：

  \begin{enumerate}
  \def\labelenumii{(\arabic{enumii})}
  \tightlist
  \item
    \([Cu(NH_3)_4]SO_4\) (2) \(K_2[PtCl_6]\) (3) \([CoCl_2(en)_2]Cl\)
  \end{enumerate}
\item
  写出下列配合物的化学式：

  \begin{enumerate}
  \def\labelenumii{(\arabic{enumii})}
  \tightlist
  \item
    六氰合铁(Ⅲ)酸钾 (2) 三氯化五氨·一水合钴(Ⅲ) (3) 四氯合铂(Ⅱ)酸四氨合铜(Ⅱ)
  \end{enumerate}
\item
  判断 \([Ni(CN)_4]^{2-}\) 和 \([Ni(NH_3)_4]^{2+}\) 的空间构型分别是什么？为什么同一中心离子配位数相同但构型不同？
\item
  解释为什么 \([Cu(NH_3)_4]SO_4\) 溶液中加 BaCl2 产生白色沉淀，加 NaOH 却不产生蓝色沉淀。
\end{enumerate}

\subsubsection{6.2 提高训练（⭐⭐⭐）}\label{ux63d0ux9ad8ux8badux7ec3}

\begin{enumerate}
\def\labelenumi{\arabic{enumi}.}
\setcounter{enumi}{4}
\item
  实验测得 \([Fe(H_2O)_6]^{2+}\) 的磁矩为 5.30 B.M.，\([Fe(CN)_6]^{4-}\) 的磁矩为 0 B.M.。请分别判断它们的 d 电子排布、杂化类型、内/外轨型和空间构型。
\item
  验证下列配合物是否符合 18 电子规则：

  \begin{enumerate}
  \def\labelenumii{(\arabic{enumii})}
  \tightlist
  \item
    \([Cr(CO)_6]\) (2) \([Mn_2(CO)_{10]}\) (3) \([Fe(CO)_4H_2]\) (4) \([V(CO)_6]\)
  \end{enumerate}
\item
  计算下列配合物在八面体场中的 CFSE（以 Dq 为单位）：

  \begin{enumerate}
  \def\labelenumii{(\arabic{enumii})}
  \tightlist
  \item
    \([Ti(H_2O)_6]^{3+}\)（d¹） (2) \([Co(NH_3)_6]^{3+}\)（低自旋 d⁶） (3) \([Mn(H_2O)_6]^{2+}\)（高自旋 d⁵）
  \end{enumerate}
\item
  画出 \([Pt(NH_3)_2Cl_2]\) 的所有几何异构体，并指出哪个有抗癌活性。
\item
  计算 \([Ni(CN)_4]^{2-}\) 中 Ni2+ 的配位数。已知 CN- 为单齿配体。
\end{enumerate}

\subsubsection{6.3 挑战题（⭐⭐⭐⭐）}\label{ux6311ux6218ux9898}

\begin{enumerate}
\def\labelenumi{\arabic{enumi}.}
\setcounter{enumi}{9}
\item
  某八面体配合物 \([Co(NH_3)_4Cl_2]^+\) 的磁矩为 0 B.M.。

  \begin{enumerate}
  \def\labelenumii{(\arabic{enumii})}
  \tightlist
  \item
    判断中心离子的 d 电子排布和自旋状态。
  \item
    推测该配合物可能的几何异构体数目，并判断各有无光学活性。
  \item
    若将 Cl- 替换为 Br-，异构体数目是否变化？为什么？
  \end{enumerate}
\item
  \(Co^{3+}\) 与 en（乙二胺）形成配合物 \([Co(en)_3]^{3+}\)。

  \begin{enumerate}
  \def\labelenumii{(\arabic{enumii})}
  \tightlist
  \item
    画出该配合物的全部异构体。
  \item
    \([Co(en)_2Cl_2]^+\) 有多少种异构体？其中哪些有光学活性？
  \item
    用晶体场理论计算 \([Co(en)_3]^{3+}\) 的 CFSE（已知 en 为强场配体，d⁶ 低自旋）。
  \end{enumerate}
\item
  已知 \([Ti(H_2O)_6]^{3+}\) 在 490 nm 处有最大吸收。计算其 Δₒ（单位 kJ/mol）并预测配合物颜色。
\item
  比较 \([CoF_6]^{3-}\) 和 \([Co(CN)_6]^{3-}\) 的自旋状态、磁矩和 CFSE，解释差异原因。
\item
  \(K_3[FeF_6]\) 的磁矩为 5.92 B.M.，而 \(K_3[Fe(CN)_6]\) 的磁矩约为 2.3 B.M.。解释磁矩差异的电子结构原因。
\item
  解释为什么 d⁵ 高自旋的 CFSE = 0，但 \([Mn(H_2O)_6]^{2+}\) 仍然可以稳定存在。
\end{enumerate}

\subsection{七、参考答案与提示}\label{ux4e03ux53c2ux8003ux7b54ux6848ux4e0eux63d0ux793a}

\textbf{6.1 基础巩固}

\begin{enumerate}
\def\labelenumi{\arabic{enumi}.}
\item
  \begin{enumerate}
  \def\labelenumii{(\arabic{enumii})}
  \tightlist
  \item
    硫酸四氨合铜(Ⅱ)，Cu2+，配位数 4
  \item
    六氯合铂(Ⅳ)酸钾，Pt⁴+，配位数 6
  \item
    一氯化二氯·二(乙二胺)合钴(Ⅲ)，Co3+，配位数 6
  \end{enumerate}
\item
  \begin{enumerate}
  \def\labelenumii{(\arabic{enumii})}
  \tightlist
  \item
    \(K_3[Fe(CN)_6]\) (2) \([Co(NH_3)_5(H_2O)]Cl_3\) (3) \([Cu(NH_3)_4][PtCl_4]\)
  \end{enumerate}
\item
  \([Ni(CN)_4]^{2-}\)：dsp2 → 平面正方形；\([Ni(NH_3)_4]^{2+}\)：sp3 → 正四面体。CN- 为强场配体，使 Ni2+ 的 3d 电子重排归并，空出 3d 轨道形成 dsp2 杂化。
\item
  \([Cu(NH_3)_4]SO_4 \rightarrow [Cu(NH_3)_4]^{2+} + SO_4^{2-}\)。SO42- 在外界，遇 Ba2+ 产生 BaSO4 白色沉淀；Cu2+ 在内界被 NH3 牢固配位，不释放 Cu2+，故无 Cu(OH)2 沉淀。
\end{enumerate}

\textbf{6.2 提高训练}

\begin{enumerate}
\def\labelenumi{\arabic{enumi}.}
\setcounter{enumi}{4}
\item
  \textbf{\([Fe(H_2O)_6]^{2+}\)}（μ = 5.30 → n = 4）→ d⁶ 高自旋 \(t_{2g^4 e_g^2}\) → sp3d2 外轨 → 正八面体
  \textbf{\([Fe(CN)_6]^{4-}\)}（μ = 0 → n = 0）→ d⁶ 低自旋 \(t_{2g^6 e_g^0}\) → d2sp3 内轨 → 正八面体
\item
  \begin{enumerate}
  \def\labelenumii{(\arabic{enumii})}
  \tightlist
  \item
    Cr(0): 6e- + 6×CO=12 → 18 ✅
  \item
    Mn(0): 7e-，每个 Mn: 5×CO端基=10 + 1×Mn---Mn键=1 → 18 ✅
  \item
    Fe(II): 6e- + 4×CO=8 + 2×H-=4 → 18 ✅
  \item
    V(0): 5e- + 6×CO=12 → 17 ❌（奇电子，空间位阻阻止二聚）
  \end{enumerate}
\item
  \begin{enumerate}
  \def\labelenumii{(\arabic{enumii})}
  \tightlist
  \item
    d¹ → \(t_{2g^1}\) → CFSE = \(-4\ Dq\)
  \item
    d⁶ 低自旋 → \(t_{2g^6}\) → CFSE = \(-24\ Dq\)
  \item
    d⁵ 高自旋 → \(t_{2g^3 e_g^2}\) → CFSE = 0
  \end{enumerate}
\item
  Pt2+ 为 d⁸，平面正方形。顺式（cis）有抗癌活性，反式（trans）无抗癌活性。
\item
  Pt2+ 配位数 = 4（4 个单齿 CN- 配体）
\end{enumerate}

\textbf{6.3 挑战题}

\begin{enumerate}
\def\labelenumi{\arabic{enumi}.}
\setcounter{enumi}{9}
\item
  \begin{enumerate}
  \def\labelenumii{(\arabic{enumii})}
  \tightlist
  \item
    μ = 0 → n = 0 → Co3+ d⁶ 低自旋 → \(t_{2g^6 e_g^0}\)
  \item
    \([Co(NH_3)_4Cl_2]^+\)（\(Ma_4e_2\) 型）→ 顺式（1,2-）和反式（1,6-）。顺式无对称面 → 有光学活性（一对对映体）；反式有对称面 → 无光学活性。总计 3 种异构体。
  \item
    替换为 Br- 后异构体数目不变，仍为 \(Ma_4e_2\) 型。
  \end{enumerate}
\item
  \begin{enumerate}
  \def\labelenumii{(\arabic{enumii})}
  \tightlist
  \item
    \([Co(en)_3]^{3+}\)（\([M(AA)_3]\) 型）→ 一对光学异构体（Λ 和 Δ）。
  \item
    \([Co(en)_2Cl_2]^+\)：反式（无光学活性）+ 顺式（一对对映体）= 共 3 种。
  \item
    en 为强场配体 → Co3+ d⁶ 低自旋 → \(t_{2g^6}\) → CFSE = \(-24\ Dq\)。
  \end{enumerate}
\item
  \(\Delta_o = N_A hc/\lambda = 6.02\times10^{23} \times 6.626\times10^{-34} \times 3.0\times10^8 / (490\times10^{-9}) \approx 239\ kJ/mol\)。吸收绿光（\textasciitilde490 nm），呈紫红色。
\item
  \([CoF_6]^{3-}\)：F- 弱场 → 高自旋 d⁶（\(t_{2g^4 e_g^2}\)，n=4，μ≈4.90 B.M.，CFSE=-4 Dq）。\([Co(CN)_6]^{3-}\)：CN- 强场 → 低自旋 d⁶（\(t_{2g^6}\)，n=0，μ=0，CFSE=-24 Dq）。Δ 与 P 的相对大小决定。
\item
  K3{[}FeF6{]}：F- 弱场 → Fe3+ d⁵ 高自旋（5 个未成对电子，μ=5.92 B.M.）。K3{[}Fe(CN)6{]}：CN- 强场 → Fe3+ d⁵ 低自旋（1 个未成对电子，μ≈2.3 B.M.，实验值因自旋-轨道耦合略高）。
\item
  d⁵ 高自旋的 \(t_{2g^3 e_g^2}\) 排布中，t2g 的 -12 Dq 与 e\_g 的 +12 Dq 恰好抵消 → CFSE = 0。但配合物的稳定性不仅来自 CFSE，还来自配位键的强度、溶剂化能、螯合效应等。\([Mn(H_2O)_6]^{2+}\) 的配位键本身就有相当的强度。
\end{enumerate}

\subsection{八、延伸阅读}\label{ux516bux5ef6ux4f38ux9605ux8bfb}

\begin{itemize}
\tightlist
\item
  配合物 --- 核心知识点详解
\item
  晶体场理论 --- 完整推导 + Jahn-Teller 效应 + CFSE 全表
\item
  18电子规则 --- 竞赛级应用详解
\item
  专题-配位化学 --- 专题知识系统总结（含解题套路+例题串讲）
\item
  12-教学洞察/教学洞察-配合物基础 --- 教学决策支持：认知冲突、高频错误、有效类比
\item
  光谱化学序列 --- 配体场强完整解析与 σ/π 解释
\item
  高自旋与低自旋 --- Δ vs P 判断逻辑 + dⁿ 组态速查
\item
  过渡金属通性 --- 配合物的元素化学特征
\item
  宋天佑《无机化学》§8-9 章（配合物 + 晶体场）
\item
  上海中学竞赛课程·化学·第二分册·第六讲
\end{itemize}

\end{document}
